\documentclass[10pt]{article}
\usepackage{kotex}
\usepackage[margin=18mm]{geometry}
\usepackage{booktabs,longtable,tabularx,array,xcolor,hyperref}
\usepackage{siunitx}
\hypersetup{hidelinks}
\setlength{\parindent}{0pt}
\setlength{\parskip}{0.45em}
\newcolumntype{Y}{>{\raggedright\arraybackslash}X}
\title{RSA-free Sahai-POC와 zkDPP-POC 성능 평가 보고서}
\author{gnark PLONK-KZG 및 EVM 재현 실험}
\date{2026년 8월}
\begin{document}
\maketitle
\tableofcontents
\newpage
\section{목적과 핵심 비교 질문}
이 보고서는 Sahai et al. (2020)의 여섯 supply-chain Event predicate를 재현한 \textbf{Sahai-POC}와 기존 \textbf{zkDPP-POC}의 성능을 비교한다. 공식 Sahai-POC는 RSA accumulator를 제거하고 \texttt{DocHash}를 commitment 겸 ledger document ID로 사용한다. 모든 input은 한 번만 소비되는 hardened semantics다.

핵심 질문은 (1) 같은 Poseidon2/PLONK/EVM 환경에서 Event별 gas와 live-proof E2E가 얼마인가, (2) SHA-256을 쓰면 constraint와 proving cost가 얼마나 변하는가, (3) Anvil과 permissioned Besu QBFT에서 EVM gas가 같은가이다. 두 protocol의 Event semantics는 완전히 같지 않으므로 기능적으로 가장 가까운 관계의 실제 값을 나란히 제시하며 동등 protocol이라고 주장하지 않는다.

\subsection{용어와 측정 원칙}
\begin{description}
\item[Sahai-POC] RSA-free, \texttt{DocHash}-ID adaptation. 직접 provenance는 \texttt{InDocs}로 유지한다.
\item[zkDPP-POC] 기존 \texttt{poc-v2} protocol/circuit. 원본 checkout은 수정하지 않고 임시 복사본에서 실행했다.
\item[MT] \texttt{GOMAXPROCS=8}. 독립 proof 호출은 순차적이고 gnark 내부 병렬성만 사용한다. 공식 Event 비교는 MT만 사용한다.
\item[ST] \texttt{GOMAXPROCS=1}. 원 논문 Table III와 gadget 단위로 비교하기 위한 진단값이며 공식 Event 비교에는 사용하지 않는다.
\item[Gas] fixed proof transaction receipt의 \texttt{gasUsed}. 통화 단위 fee가 아니며 proof 생성 시간과 별개다.
\item[E2E] build, witness, 순차 prove, ABI encode, transaction 제출부터 receipt까지의 시간. circuit compile/setup, artifact load, deployment는 제외한다.
\end{description}
각 공식 조합은 한 번만 실행했다. 같은 Event가 시나리오에 여러 번 나타나면 독립 반복 run이 아니라 서로 다른 case이며, 표에는 case 수와 case 중앙값을 쓴다.
\section{연결 canonical 시나리오}
\begin{longtable}{cllllrr}
\toprule 순서 & Case & Event & Sender & Recipient & 입력 수량 & 출력 수량\\
\midrule\endhead
1 & entry\_a & Entry & Supplier A & Supplier A & -- & 100\\
2 & ship\_a & Ship & Supplier A & Processor & 100 & 100\\
3 & split\_a & Split & Processor & Processor & 100 & 40+60\\
4 & merge\_a & Merge & Processor & Manufacturer & 40+60 & 100\\
5 & entry\_b & Entry & Supplier B & Supplier B & -- & 50\\
6 & ship\_b & Ship & Supplier B & Manufacturer & 50 & 50\\
7 & process\_product & Process & Manufacturer & Exit actor & 100+50 & 150\\
8 & exit\_product & Exit & Exit actor & Exit actor & 150 & 150 terminal\\
\bottomrule
\end{longtable}
이 표는 공식 성능 비교에서 사용하는 연결된 8-transaction DAG를 정의한다. 각 output \texttt{DocHash}는 다음 Event의 \texttt{inputId}와 \texttt{InDocs}에 실제로 연결된다. Entry와 Ship은 각각 두 case, 나머지는 한 case다. Exit은 새 terminal document를 만들고 output의 \texttt{Sender=Recipient}를 Path 1개와 Eq 1개로 증명한 뒤 input을 consumed 처리한다. zkDPP-POC는 자체 battery/material 시나리오를 사용하므로 case 수와 predicate가 다르다.
\section{실험 환경과 측정 경계}
\begin{tabularx}{\textwidth}{lY}
\toprule 항목 & 값\\\midrule
CPU & Apple M1 Pro (10 logical cores)\\
Go / gnark & go1.25.7 / v0.15.0, gnark-crypto v0.20.1\\
Proof & PLONK-KZG/BLS12-381, Poseidon2 기본 / SHA-256 ablation\\
EVM & Foundry 1.7.1, Solidity 0.8.30, Prague, chain ID 31337, block gas 30M\\
Besu & 26.7.1, 4 validators + 1 RPC, QBFT block period 1초\\
실행 횟수 & profile/thread/backend의 공식 조합당 1회\\
\bottomrule
\end{tabularx}
Gas는 fixed proof의 top-level receipt만 포함하며 deployment와 setup용 state seed는 제외한다. E2E는 fresh proof 생성과 submit-to-receipt를 포함한다. MT에서도 Event 안 proof를 동시에 만들지 않으므로 protocol-level parallel batching 효과는 포함되지 않는다.
\section{Poseidon2와 SHA-256 gadget 비교}
\begin{longtable}{llrrrrrr}
\toprule Profile & Gadget & Constraint & Public & Proof B & ST prove ms & MT prove ms & MT verify ms\\
\midrule\endhead
Poseidon2 & MerklePath & 6582 & 5 & 1056 & 953.8 & 152.4 & 2.738\\
Poseidon2 & Eq & 2231 & 2 & 1056 & 517.6 & 83.2 & 2.596\\
Poseidon2 & Add & 3347 & 3 & 1056 & 556.8 & 89.5 & 2.635\\
Poseidon2 & And & 2582 & 3 & 1056 & 543.5 & 88.9 & 2.602\\
SHA-256 & MerklePath & 1071658 & 10 & 1184 & 149863.1 & 24054.7 & 2.757\\
SHA-256 & Eq & 561113 & 4 & 1184 & 75928.9 & 11231.9 & 2.760\\
SHA-256 & Add & 644678 & 6 & 1184 & 76506.2 & 11372.2 & 2.838\\
SHA-256 & And & 644681 & 6 & 1184 & 77932.1 & 11545.2 & 2.722\\
\bottomrule\end{longtable}
이 표는 hash profile이 독립 PLONK gadget의 constraint와 prove/verify에 미치는 영향을 묻는다. 각 행은 동일 fixture와 key로 ST 1회, MT 1회 실행한 값이다. verify는 Go native verifier 1회이며 Solidity gas가 아니다. SHA-256은 bit-level compression 때문에 Poseidon2보다 constraint와 prove 시간이 크게 증가하지만 proof byte 차이는 작다. ST는 여기와 마지막 원 논문 비교에서만 사용한다.
\section{Sahai-POC Event payload와 hash ablation}
\begin{longtable}{lrrrrrrr}
\toprule Event & Cases & Proof & P2 logical B & SHA logical B & P2 prove ms & SHA prove ms & SHA/P2\\
\midrule\endhead
Entry & 2 & 2 & 5207 & 5742 & 227.9 & 32671.6 & 143.3\\
Ship & 2 & 5 & 12582 & 13915 & 619.9 & 79260.1 & 127.9\\
Merge & 1 & 6 & 15413 & 17027 & 678.6 & 98314.8 & 144.9\\
Split & 1 & 6 & 15343 & 16953 & 679.9 & 98313.2 & 144.6\\
Process & 1 & 6 & 15425 & 17043 & 779.8 & 98497.4 & 126.3\\
Exit & 1 & 2 & 5390 & 5922 & 239.1 & 32941.8 & 137.8\\
\bottomrule\end{longtable}
이 표는 연결 시나리오에서 Event bundle을 만들 때 hash 선택이 payload와 순차 proof 생성에 미치는 영향을 보여준다. Logical byte는 JSON payload이며 ABI calldata와 다르다. Entry/Ship은 두 case 중앙값, 나머지는 한 case 값이다. Proof 수는 각각 Entry 2, Ship 5, Merge/Split/Process 6, Exit 2다. 모두 MT이며 setup과 artifact load는 제외한다. 마지막 열은 SHA-256 prove / Poseidon2 prove다.
\subsection{Anvil에서의 hash profile 영향}
\begin{longtable}{lrrrrrr}
\toprule Event & Cases & P2 gas & SHA gas & Gas 비 & P2 E2E ms & SHA E2E ms\\
\midrule\endhead
Entry & 2 & 809029 & 862947 & 1.067 & 242.1 & 36613.9\\
Ship & 2 & 1917717 & 2050816 & 1.069 & 567.4 & 89730.2\\
Merge & 1 & 2308199 & 2471327 & 1.071 & 736.8 & 124857.2\\
Split & 1 & 2390246 & 2553386 & 1.068 & 733.7 & 111938.5\\
Process & 1 & 2308110 & 2471238 & 1.071 & 759.5 & 139062.6\\
Exit & 1 & 857814 & 911768 & 1.063 & 241.2 & 37009.2\\
\bottomrule\end{longtable}
Gas 열은 profile별 fixed proof transaction을 한 번 실행한 receipt \texttt{gasUsed}의 case 중앙값이고, E2E는 fresh proof MT 1회의 case 중앙값이다. Gas 비는 SHA-256/Poseidon2다. SHA-256의 E2E 증가는 대부분 prove에 기인하며 gas 증가는 proof/public-input calldata와 verifier 차이만 반영한다. 즉 native proof 생성비와 EVM 검증비는 서로 다른 병목이다.
\section{Anvil 기반 zkDPP-POC 비교}
\subsection{Event receipt gas}
\begin{longtable}{llrrrrr}
\toprule Sahai & zkDPP 대응 & S cases & Z cases & Sahai gas & zkDPP gas & S/Z\\
\midrule\endhead
Entry & entry & 2 & 4 & 809029 & 1498540 & 0.540\\
Ship & transfer & 2 & 7 & 1917717 & 2673540 & 0.717\\
Merge & merge & 1 & 1 & 2308199 & 1561882 & 1.478\\
Split & split & 1 & 1 & 2390246 & 2352560 & 1.016\\
Process & process & 1 & 3 & 2308110 & 2403091 & 0.960\\
Exit & exit & 1 & 2 & 857814 & 401760 & 2.135\\
\bottomrule\end{longtable}
이 표는 Poseidon2 MT의 실제 gas를 먼저 보여주고 마지막 열에서만 Sahai-POC/zkDPP-POC 비율을 계산한다. 각 protocol의 연결 시나리오 case 중앙값이며 조합당 fixed-proof run은 1회다. Entry와 Ship은 Sahai-POC의 gas가 더 작지만 Merge와 Exit은 더 크다. 이는 proof 수만이 아니라 각 protocol의 verifier 호출, calldata, state transition 차이를 함께 반영한다.

\subsection{MT live-proof E2E}
\begin{longtable}{llrrrrr}
\toprule Sahai & zkDPP 대응 & S cases & Z cases & Sahai ms & zkDPP ms & S/Z\\
\midrule\endhead
Entry & entry & 2 & 4 & 242.1 & 98.8 & 2.449\\
Ship & transfer & 2 & 7 & 567.4 & 492.8 & 1.151\\
Merge & merge & 1 & 1 & 736.8 & 482.7 & 1.526\\
Split & split & 1 & 1 & 733.7 & 495.5 & 1.481\\
Process & process & 1 & 3 & 759.5 & 884.8 & 0.858\\
Exit & exit & 1 & 2 & 241.2 & 272.4 & 0.885\\
\bottomrule\end{longtable}
이 표는 양쪽 모두 Poseidon2, MT, proof 호출 순차 실행 조건의 E2E를 비교한다. 실제 값을 먼저 제시하고 마지막 열에서 Sahai-POC/zkDPP-POC 비율을 계산한다. Process와 Exit에서는 Sahai-POC가 빠르지만 나머지에서는 느리다. \texttt{proceed}와 \texttt{recall}은 Sahai에 대응 Event가 없어 제외했다. 같은 이름의 Event도 commitment, nullifier/ownership, predicate와 scenario가 달라 절대적인 protocol 우열로 해석할 수 없다.
\section{Besu QBFT 결과}
고정 MT proof 파일의 checksum과 calldata가 Anvil/Besu에서 같을 때 여섯 Event와 두 baseline의 \texttt{gasUsed}가 모두 정확히 일치했다. account allowlist 밖 transaction, allowlist 밖 node 연결, Contract participant role 없는 호출은 모두 거부되었다. validator 4개 중 1개를 중단한 상태에서도 transaction이 finality를 얻었다.
\begin{longtable}{lrrrr}
\toprule Event & Cases & Anvil receipt ms & Besu receipt ms & Besu E2E ms\\
\midrule\endhead
Entry & 2 & 14.4 & 1028.4 & 1269.0\\
Ship & 2 & 27.4 & 1022.9 & 1559.5\\
Merge & 1 & 42.3 & 948.8 & 1655.7\\
Split & 1 & 40.6 & 1115.8 & 1823.9\\
Process & 1 & 36.3 & 985.5 & 1706.4\\
Exit & 1 & 13.8 & 879.9 & 1104.3\\
\bottomrule\end{longtable}
이 표는 Poseidon2 MT live-proof 1회에서 backend finality latency를 분리해 보여준다. Besu E2E는 prove와 QBFT submit-to-receipt를 함께 포함한다. Anvil은 즉시 local mining이므로 receipt가 수십 ms인 반면 Besu는 약 1초 block period의 영향을 받는다. validator 1개 중단 Gate의 별도 transaction finality는 3019.3 ms였다. 이 section은 permissioned backend overhead를 설명하며 zkDPP 비교 비율에는 사용하지 않는다.
\section{RSA-free adaptation으로 제외된 기능과 해석 한계}
제거한 항목은 \texttt{DocPrime}, 500-bit prime certificate, \texttt{DocAccumulator}, membership/non-membership proof, RPoKE와 EVM MODEXP 검증이다. 따라서 \texttt{InDocs}를 따라 직접 provenance DAG를 조회할 수는 있지만 accumulator로 upstream 전체 set을 압축하거나 contamination membership/non-membership을 검증하는 기능은 없다. 같은 document 내용도 salt가 다르면 다른 \texttt{DocHash}가 되며 contract는 중복 \texttt{DocHash}를 거부한다.

이 제거는 PLONK gadget의 constraint를 줄이지 않는다. RSA는 gnark circuit 밖 component였기 때문이다. 대신 payload, calldata, native RSA 시간과 RSA-on-EVM gas가 사라진다. 결과는 “원 Sahai protocol 전체”가 아니라 Event predicate workload의 RSA-free adaptation이며 production 보안 또는 정확한 2020 hardware 시간 재현을 주장하지 않는다.
\section{Sahai 원 논문과의 비교}
이 section을 마지막에 둔 이유는 본 실험의 핵심 질문이 zkDPP-POC와의 비교이기 때문이다. 논문 Table II의 RSA component와 Table IV의 non-membership은 공식 구현에서 제거되었으므로 stale POC 숫자를 제시하지 않는다.
\subsection{논문 Table III와 SHA-256/Poseidon2 gadget}
\paragraph{Circuit 규모.}
{\small\setlength{\tabcolsep}{5pt}
\begin{longtable}{lcrrrr}
\toprule & \multicolumn{1}{c}{원 논문 SHA-256} & \multicolumn{2}{c}{Sahai-POC SHA-256} & \multicolumn{2}{c}{Sahai-POC Poseidon2}\\
\cmidrule(lr){2-2}\cmidrule(lr){3-4}\cmidrule(lr){5-6}
Gadget & Constraints & Constraints & Public inputs & Constraints & Public inputs\\
\midrule\endhead
MerklePath & 미보고 & 1071658 & 10 & 6582 & 5\\
G-Add & 미보고 & 644678 & 6 & 3347 & 3\\
G-Eq & 미보고 & 561113 & 4 & 2231 & 2\\
\bottomrule\end{longtable}
}
이 표는 각 gadget의 circuit 규모를 보여준다. Constraints는 PLONK circuit이 증명해야 할 제약식 수이고, Public inputs는 verifier에게 공개되는 field element 수다. 원 논문은 libsnark 구현의 constraint와 public input 수를 보고하지 않았으므로 임의로 추정하지 않는다. SHA-256은 bit-level hash 연산으로 Poseidon2보다 약 163--251배 많은 constraint를 사용하며, 이 차이가 아래 prove 시간 차이의 주요 원인이다. Public inputs은 SHA-256 digest를 128-bit limb 2개로 표현하기 때문에 Poseidon2보다 2배 많다.

\paragraph{Proof 생성 시간.}
{\small\setlength{\tabcolsep}{3pt}
\begin{longtable}{lrrrrrr}
\toprule & \multicolumn{2}{c}{원 논문 SHA-256} & \multicolumn{2}{c}{Sahai-POC SHA-256} & \multicolumn{2}{c}{Sahai-POC Poseidon2}\\
\cmidrule(lr){2-3}\cmidrule(lr){4-5}\cmidrule(lr){6-7}
Gadget & ST & MT & ST & MT & ST & MT\\
\midrule\endhead
MerklePath & 16000.0 & 5200.0 & 149863.1 & 24054.7 & 953.8 & 152.4\\
G-Add & 4950.0 & 1600.0 & 76506.2 & 11372.2 & 556.8 & 89.5\\
G-Eq & 3200.0 & 1000.0 & 75928.9 & 11231.9 & 517.6 & 83.2\\
\bottomrule\end{longtable}
}
두 Sahai-POC profile은 각각 ST 1회(\texttt{GOMAXPROCS=1})와 MT 1회(\texttt{GOMAXPROCS=8})를 측정했다. 단위는 ms이며 작을수록 빠르다. Sahai-POC SHA-256은 원 논문의 hash workload와 가장 가깝지만, libsnark 기반 원 논문과 gnark PLONK-KZG/BLS12-381 기반 POC의 proof system·curve·hardware가 다르므로 직접적인 속도 우열 비교는 아니다. Poseidon2열은 원 논문 재현값이 아니라 본 보고서의 기본 profile에서 hash-friendly circuit을 사용했을 때의 ablation 결과다.

SHA-256 circuit은 MerklePath 1,071,658개, G-Add 644,678개, G-Eq 561,113개 constraint를 사용하지만 Poseidon2는 각각 6,582개, 3,347개, 2,231개다. 따라서 Sahai-POC SHA-256의 느린 prove는 Eq/Add 연산 자체보다 bit-level SHA-256을 PLONK circuit에서 재계산하는 비용에서 주로 발생한다. 반면 MT는 ST 대비 약 6배 빠르므로, 9.1의 차이를 thread 활용 실패로 해석하지 않는다.

\paragraph{Proof 검증 시간.}
{\small\setlength{\tabcolsep}{3pt}
\begin{longtable}{lrrrrrr}
\toprule & \multicolumn{2}{c}{원 논문 SHA-256} & \multicolumn{2}{c}{Sahai-POC SHA-256} & \multicolumn{2}{c}{Sahai-POC Poseidon2}\\
\cmidrule(lr){2-3}\cmidrule(lr){4-5}\cmidrule(lr){6-7}
Gadget & ST & MT & ST & MT & ST & MT\\
\midrule\endhead
MerklePath & 35.000 & 48.000 & 3.891 & 2.757 & 3.557 & 2.738\\
G-Add & 35.000 & 48.000 & 3.680 & 2.838 & 3.567 & 2.635\\
G-Eq & 35.000 & 48.000 & 3.572 & 2.760 & 3.550 & 2.596\\
\bottomrule\end{longtable}
}
단위는 ms이며 Sahai-POC의 검증은 Go native verifier 1회 결과이다. Solidity verifier gas와 transaction receipt 시간은 포함하지 않는다. 원 논문은 proving에서 MT가 ST보다 빠르지만, 짧은 verification은 35 ms에서 48 ms로 증가한다. 논문은 이 역전의 원인을 별도로 분석하지 않으며, Sahai-POC의 one-shot native verify와도 직접 비교하지 않는다.

현재 표의 Sahai-POC 값은 \texttt{LeafIndex}를 public statement로 binding하기 전 artifact에서 측정한 값이다. Protocol conformance 수정 후 circuit·key·verifier를 재생성하고 공식 수치를 다시 측정한다.

\newpage
\subsection{논문 Table V와 RSA-free payload}
아래 표는 논문 payload와 RSA-free POC payload의 구성 차이를 확인한다. POC 열은 SHA-256 MT 연결 시나리오의 case 중앙값이다. 논문 byte에는 RSA/prime/accumulator 관련 구성이 포함되지만 POC는 이를 제거했고 JSON logical payload를 측정했으므로 byte 값은 직접 동등 비교가 아니다. Proof 개수는 Event predicate bundle을 대응시키는 확인용이며 Split/CombineAcc 표 불일치는 RSA-free 범위 밖이다.
\begin{center}
\begin{tabular}{lrrrrr}
\toprule Event & 논문 proof & 논문 byte & 논문 MT s & POC proof & POC logical byte\\
\midrule
Entry & 2 & 640 & 6.2 & 2 & 5742\\
Ship & 5 & 1216 & 13.4 & 5 & 13915\\
Merge & 6 & 1404 & 19.2 & 6 & 17027\\
Split & 6 & 1600 & 19.2 & 6 & 16953\\
Process & 6 & 1404 & 19.2 & 6 & 17043\\
Exit & 2 & 640 & 6.3 & 2 & 5922\\
\bottomrule\end{tabular}
\end{center}
\end{document}
